\documentclass[10pt,letterpaper]{article}

\usepackage{amssymb}
\usepackage[T1]{fontenc}
\usepackage{fullpage}
\usepackage{hyperref}
\usepackage{lmodern}
\usepackage{needspace}
\usepackage{tabularx}
\usepackage{xcolor}
\usepackage{xspace}
\usepackage{tikz}			% custom graphics

% Header and section style follow the VPK research statement template.
\def\name{\Huge{Vaggelis Atlidakis}\xspace}
\def\type{\LARGE{\textsl{Research Statement}}\xspace}
\def\ie{\emph{i.e.,}\xspace}

\hypersetup{
  colorlinks = true,
  citecolor = black,
  filecolor = black,
  linkcolor = black,
  urlcolor = black,
  pdfauthor = {Vaggelis Atlidakis},
  pdftitle = {Research Statement}
}

\definecolor{green}{rgb}{0.55,0.85,0.35}
\newcommand*{\colorrule}[2][1ex]{{\color{green}{\rule{#2}{#1}}}}
\renewcommand{\thesection}{\colorrule[0.7ex]{.125\linewidth}}
% Keep each project heading with enough of its description to read in context.
\newcommand{\project}[2]{\par\Needspace{6\baselineskip}\noindent
  $\blacktriangleright$~\textbf{#1~\cite{#2}.}}
\widowpenalty=10000
\clubpenalty=10000

\begin{document}
\fontsize{10.5pt}{12.4pt}\selectfont

{\centering
\begin{tabularx}{\textwidth}{X r}
  \name & \type \\
\end{tabularx}
{\color{green}\rule{\linewidth}{.25ex}}}
\par\vspace{0.45\baselineskip}

% circled numbers
\newcommand*\circled[1]{\tikz[baseline=(char.base)]{
	\node[shape=circle,draw,inner sep=0.75pt] (char) {#1};}}

\noindent
\emph{My research focuses on building secure and reliable software systems.}
The same complexity that makes modern software useful also hides failures:
cloud APIs require long sequences of operations, machine-learning frameworks
cross language and hardware boundaries, and operating systems expose broad
interfaces to untrusted code. AI coding agents will create and modify software
at still greater speed. My research addresses the resulting security and
reliability challenge in three connected ways: \emph{discovering errors through
automated testing}, \emph{limiting their consequences through systems
hardening}, and \emph{prioritizing remediation using evidence from running
software}. I build systems that turn these ideas into deployable tools and
measure them against real applications.

My research method begins with the behavior of real software. An API
specification does not describe every state a cloud service can enter; a
high-level ML interface does not reveal which native or GPU kernel will run;
and a standard OS interface may not express the abstractions an application
actually needs. I measure these gaps, then use the structure I uncover to
build tests that reach otherwise hidden behavior. RESTler uses request
dependencies and response feedback, IvySyn and FuzzWarp cross language and
device boundaries, and my POSIX study measures how applications use OS
interfaces in practice. The common goal is to produce evidence that a
developer can reproduce and act on.

Discovery also informs protection. EPF exposes how a kernel interface can
undermine isolation, while SysXCHG, FineIBT, and Eclipse narrow privileges or
block exploit paths without requiring every bug to be removed first. At
Lacework, I applied the same empirical approach to remediation: runtime
package activity helps teams decide which known vulnerabilities need
attention now. Across these projects, I seek mechanisms that can be
implemented in existing systems, measured for cost and effectiveness, and
used beyond a research prototype.

This agenda has developed across academic and industrial settings. My
doctoral work at Columbia and Microsoft Research taught me to make testing
follow the state of a real cloud service rather than isolated API calls. At
Brown, I extended that attention to execution context across language,
device, kernel, and hardware boundaries, and worked on defenses that remain
practical in existing software stacks. At Lacework, I saw how deployment
changes the question from whether a vulnerability exists to which finding a
team should address first. Together, these experiences have shaped my aim to
make security research useful at each stage, from discovering a failure to
deploying a fix.

{\noindent$\blacktriangleright$~\textbf{Impact and Technology Transfer.}}
\textbf{RESTler}~\cite{restler,restler2}, the first stateful REST API fuzzer,
has been used by Microsoft teams since 2017 for security, reliability, and
regression testing. It found hundreds of bugs in production O365, Azure, and
Bing services, and Microsoft open-sourced it in 2020.
\textbf{FineIBT}~\cite{fineibt} was adopted by Intel, and a variant was
upstreamed to the Linux kernel. At Lacework, I helped develop
\textbf{Code-Aware Agent (CAA)}~\cite{caa}, which brings continuous package
activity monitoring into production vulnerability prioritization. A reported
study of 2,428 machines found that 79\% of vulnerable packages were inactive
during its observation period, showing the practical value of runtime context.

The tools also prompted concrete fixes in widely used software.
\textbf{IvySyn}~\cite{ivysyn} helped TensorFlow and PyTorch developers
identify and fix 61 previously unknown vulnerabilities, resulting in 39 CVEs.
\textbf{FuzzWarp}~\cite{fuzzwarp} found 99 previously unknown GPU-specific
memory-safety bugs in those frameworks; at the time of publication, 62 had
been acknowledged and 37 had been fixed or had approved fixes.
\textbf{Cynthia}~\cite{cynthia} found 28 bugs in five object-relational
mapping systems; 25 were confirmed, 20 were fixed, and three were marked as
release blockers. These results illustrate why I pair automated discovery
with reproducible reports and direct engagement with maintainers.
\pagebreak

\section{Automated Software Testing}
XXX My work on automated software testing follows a model-based principle:
RESTler~\cite{restler,restler2}
IvySyn~\cite{ivysyn}
FuzzWarp~\cite{fuzzwarp}
StratoFuzz~\cite{stratofuzz}
FairTest~\cite{fairtest}
Cynthia~\cite{cynthia}
Pythia~\cite{pythia}
effective tests must capture the state, dependencies, execution context, or
expected behavior that shape a system's response. Specifications, feedback,
and executable models guide valid sequences through cloud APIs and satellite
p rotocols; control of execution paths brings tests to native and GPU kernels;
p rotocols; control of execution paths brings tests to native and GPU kernels;
p rotocols; control of execution paths brings tests to native and GPU kernels;
p rotocols; control of execution paths brings tests to native and GPU kernels;
p rotocols; control of execution paths brings tests to native and GPU kernels;
p rotocols; control of execution paths brings tests to native and GPU kernels;
and abstract queries, targeted data, and statistical models reveal silent
correctness and fairness failures. Across these settings, I use system
structure to make exploration tractable and produce reproducible evidence
that developers can act on.
XXX

\project{RESTler}{restler}
% OK: DO NOT TOUCH
During my Ph.D., along with collaborators at Microsoft Research, we introduced
stateful REST API fuzzing and built RESTler, the first stateful REST API fuzzer.
Although REST APIs expose cloud-service operations individually, our key
observation was that those
operations are not independent: their behavior can depend on resources created
by earlier requests. A fuzzer must, therefore, generate request sequences that
establish the state needed to reach failures missed by isolated tests. RESTler
analyzes an OpenAPI specification to infer producer-consumer dependencies:
for example, a request that creates a project produces an identifier needed by
later requests that create branches or issues within that project. It, then,
combines this static structure with dynamic feedback from service responses,
reusing resource identifiers and pruning sequences that the service rejects.
Our evaluation showed that both dependency inference and response feedback were
necessary to explore deeply while controlling the combinatorial search space.
This work
established a way to automatically test cloud services through their APIs.
Since 2017, RESTler is used internally by multiple Microsoft
teams for security, reliability, and functional correctness testing,
and has found hundreds of bugs in multiple production O365, Azure, and Bing
cloud services.
In 2020, RESTler became open-source, and its usage has been
steadily growing since, both inside and outside Microsoft.

\project{Checking security properties of REST APIs}{restler2}
% OK: DO NOT TOUCH
Crashes exposed by RESTler cover only part of a cloud service's attack surface:
an API may return a
successful response while violating resource ownership or isolation. Therefore,
we extended RESTler with four security rules covering access to deleted
resources, side effects of failed resource creation, parent-child resource
isolation, and
separation between user namespaces. Then, we implemented modular
active checkers that generate additional requests specifically intended to
violate these rules.
The checkers integrate with stateful exploration, which
first brings the service into the state needed to test a property.
Across nearly a dozen production Azure and Office 365 services tested, about
one third of
the bugs found were rule violations rather than internal-server errors---showing
that automated API testing can
detect non-crashing security failures with implications for resource hijacking, information
disclosure, and denial of service.

\project{Pythia}{pythia}
% OK: DO NOT TOUCH
RESTler's generated grammars preserve request structure and dependencies, but
inherently use fixed fuzzing dictionaries of strings and numbers, which limit
the set of malformed value combinations that can be explore. To address this,
I co-developed Pythia, which uses a statistical model to learn common usage
patterns of a target service REST API from structurally valid seed inputs, and
then generates learning-based mutations by injecting a small amount of noise
deviating from common usage patterns, while still maintaining syntactic validity.
Pythia's mutation strategy helps generate grammatically valid test cases, and
coverage-guided feedback helps prioritize the test cases that are more likely
to find bugs. 
Experimental evaluation on GitLab, Mastodon, and Spree, showed that
Pythia increased coverage after baseline stateful fuzzing had plateaued, and
also found $29$ new bugs. This work demonstrated how
to use learning to expand fuzzing domains, while preserving the structure
needed to reach deeper application logic.

\project{IvySyn}{ivysyn}
% OK: DO NOT TOUCH
Deep-learning frameworks expose native C/C++ kernels through managed-language
interfaces, and missing checks or inconsistent assumptions between the layers of
native and managed languages can let untrusted inputs reach memory-unsafe code.
To address this, along with collaborators from Brown university, we
introduced a two-fold bottom-up testing approach, implemented in
our IvySyn framework. IvySyn leverages the statically-typed nature of native
APIs in order to automatically perform type-aware mutation-based fuzzing
on low-level kernel code. Given a set of offending inputs that trigger memory
safety (and runtime) errors in low-level, native DL (C/C++) code, IvySyn
automatically synthesizes code snippets in high-level languages (e.g., in
Python), which propagate error-triggering input via high(er)-level APIs.
Such code snippets essentially act as ``Proof of Vulnerability'',
as they demonstrate the existence of bugs in native code that an attacker
can target through various high-level APIs. Our evaluation showed that IvySyn
significantly outperforms past approaches, both in terms of efficiency and
effectiveness, in finding vulnerabilities in popular DL frameworks. We used
IvySyn to test TensorFlow and PyTorch, and helped the TensorFlow and PyTorch
framework developers to identify and fix 61 previously-unknown security
vulnerabilities ($39$ unique CVEs assigned), implicitly and explicitly
affecting millions, if not billions, of users worldwide.

\project{FuzzWarp}{fuzzwarp}
% OK: DO NOT TOUCH
Our recent FuzzWarp work extends IvySyn's bottom-up testing approach to target
GPU kernels, while directly controlling the host-to-GPU execution path. FuzzWarp
first verifies that each test input reaches the intended target kernel, which was
a critical shortcoming of all prior work: in our baseline campaigns, 29--44\%
of test inputs generated from existing fuzzers did not reach the target GPU kernels.
To address this, FuzzWarp allocates each mutated
tensor through the framework's allocator on the device the kernel expects,
then synchronizes the GPU after every mutation, and uses NVIDIA's Compute
Sanitizer to attribute asynchronous memory faults to their triggering inputs.
FuzzWarp employs mutations that target dimension-product overflows, boundary values, and
zero-dimension tensors, and akin to IvySyn, it also lifts low-level
failures into reproducible Python API proof-of-vulnerability snippets. Across
multiple TensorFlow versions, FuzzWarp consistently outperformed state-of-the-art
deep-learning fuzzers
and uncovered $99$ previously unknown GPU-specific memory-safety bugs in
TensorFlow and PyTorch ($62$ confirmed; $37$ fixed).

\project{StratoFuzz}{stratofuzz}
% OK: DO NOT TOUCH
Space missions depend on communication stacks that
preserve the confidentiality, integrity, and availability of their
data links. Testing these stacks remains difficult, however,
because message validity depends on hidden state and inter-
layer temporal dependencies; while reaching deeper states
requires longer test sequences, whose number grows expo-
nentially with length. Along with collaborators from TU Delft,
we introduced StratoFuzz, a framework for black-box testing
of complete, stateful protocol stacks for satellite communication. Our key
observation is that a single executable, model of an entire such stack, can
serve as (1) a message generator, (2) a correctness oracle, and (3) a proxy for
the recipient' s state. Compared against AFLNet and StateAFL, StratoFuzz
achieved significantly higher line coverage than AFLNet and StateAFL. Neither
baseline reported a finding because their mutations failed to preserve the
inter-layer temporal dependencies needed to reach deeper protocol states.
By contrast, StratoFuzz found 29 previously unknown bugs (23 confirmed; one CVE
assigned), most requiring complex multi-message sequences in order to be
triggered

\project{Cynthia}{cynthia}
% OK: DO NOT TOUCH
ORMs can silently return incorrect query results, yet their distinct APIs and
dependence on database contents make such bugs difficult to test systematically.
With collaborators from the Athens University of Economics and Business, we
introduced what is, to the best of our knowledge, the first approach for
systematically testing ORM systems, implemented as Cynthia. Differential
testing establishes a test oracle for ORM-specific bugs. The approach first
generates random relational database schemas and sets up the corresponding
databases, then queries them through the APIs of the ORMs under test. Because
ORMs lack a common input language, it generates queries in an abstract query
language and translates them into concrete, executable ORM queries. A solver
produces database records targeted to each query's constraints, making the
comparisons informative. Cynthia found 28 bugs in five popular ORM systems:
$25$ were confirmed, of which, $20$ were fixed, and three were marked as
release blockers by their corresponding developers.

\project{FairTest}{fairtest}
% OK: DO NOT TOUCH
Data-driven applications can treat users unfairly or offensively, yet such
effects are difficult to identify without suitable tools. With collaborators
from EPFL, we introduced the unwarranted associations (UA) framework, which
unifies prior approaches to algorithmic fairness. UA combines investigative
primitives and fairness metrics with subgroup analysis and utility factors
that may explain disparities. We instantiated it in FairTest, the first
comprehensive tool for checking unfair user treatment in data-driven
applications. FairTest supports scalable, statistically rigorous analysis of
associations between outcomes and sensitive attributes, with debugging tools
to rule out confounders. Across four applications, it revealed disparate
impact, offensive labeling, and uneven algorithmic error rates, including
bias against older users in a predictive health application and offensive
racial labels in an image tagging application.

\section{Systematic Software Hardening}
% NOT-OK: XXX
XXX Testing will never find every defect. My hardening work therefore narrows the
authority available to compromised code and blocks exploit paths even when a
memory error remains. With collaborators at Brown, I have used compiler
instrumentation and existing OS and processor mechanisms to make these defenses
practical for deployed software.

\project{SysXCHG}{sysxchg}
% NOT-OK: XXX
Linux programs inherit access to a broad system-call interface, giving
compromised code a rich set of operations for further exploitation.
Conventional seccomp-BPF filtering introduces another source of over-privilege:
because filters are inherited and can only become stricter, a program that
executes other programs must allow the calls needed by its descendants.
My collaborators and I designed SysXCHG around an exchange model. Each binary
carries an \emph{exec filter}, automatically installed when that binary is
executed, so protection adapts to the current program. Explicitly installed
manual filters remain irreversible, preserving developer-imposed restrictions.
We use Linux's Integrity Measurement Architecture to sign filter-carrying
binaries and prevent tampering that would widen privileges. We also developed
\texttt{xfilter}, which enforces number-based policies through a process-specific
view of the kernel's syscall table. In the PaSH macrobenchmarks, the combined
design incurred at most 1.71\% slowdown; \texttt{xfilter} also reduced filter
installation time by approximately 76--96\%. SysXCHG therefore removes the
over-privilege caused by filter inheritance while retaining efficient
enforcement.

% We present the design, implementation, and evaluation of SysXCHG:
% a system call (syscall) filtering enforcement mechanism that enables
% programs to run in accordance with the principle of least privilege.
% In contrast to the current, hierarchical design of seccomp-BPF,
% which does not allow a program to run with a different set of allowed
% syscalls than its descendants, SysXCHG enables applications to run
% with “tight” syscall filters, uninfluenced by any future-executed
% (sub-)programs, by allowing filters to be dynamically exchanged at
% runtime during execve[at]. As a part of SysXCHG, we also present
% xfilter: a mechanism for fast filtering using a process-specific
% view of the kernel’s syscall table where filtering is performed. In
% our evaluation of SysXCHG, we found that our filter exchanging
% design is performant, incurring ≤ 1.71% slowdown on real-world
% programs in the PaSH benchmark suite, as well as effective, blocking
% vast amounts of extraneous functionality, including security-critical
% syscalls, which the current design of seccomp-BPF is unable to

% Lastly, to allow programs that invoke other programs
% (e.g., via execve[at]) execute in accordance to the PoLP, I designed SysXCHG: an innovative least-privilege
% syscall-filtering enforcement mechanism that addresses the limitations of seccomp-BPF. Specifically, in contrast to the current, hierarchical design of seccomp-BPF, which does not permit a program to run with a
% different set of allowed syscalls than its descendants, SysXCHG enables applications to run with “tight”
% syscall filters, uninfluenced by any future-executed (sub-)programs, by allowing filters to be dynamically
% exchanged at runtime (safely) during execve[at]. As a part of SysXCHG, we also introduced xfilter: a
% new mechanism for fast syscall filtering, which relies on process-specific views of the kernel’s syscall table.

\project{FineIBT}{fineibt}
% NOT-OK: XXX
Intel's Indirect Branch Tracking (IBT) restricts indirect calls and jumps to
marked entry points, but does not distinguish which marked target is
appropriate for a particular transfer. Supporting shared libraries makes this
policy especially permissive, because exported functions must also be marked.
With collaborators at Brown and Intel, I co-developed FineIBT to enforce
finer-grained control-flow integrity (CFI) on top of this hardware mechanism.
Compact instrumentation passes a policy identifier at an indirect call and
checks it at the target, preventing transfers to functions outside the allowed
class. The enforcement mechanism is independent of how the policy is derived,
so different analyses and CFI granularities can use the same instrumentation.
We implemented FineIBT in LLVM on Linux and evaluated both performance and
compatibility with modern software features. Slowdowns were at most 1.94\%
on SPEC CPU2017 and 1.92\% on the real-world applications tested, outperforming
Clang-CFI. FineIBT demonstrates how modest software instrumentation can turn
a coarse hardware primitive into a more precise, deployable defense.

% We present the design, implementation, and evaluation of FineIBT:
% a CFI enforcement mechanism that improves the precision of hardware-assisted
% CFI solutions, like Intel IBT, by instrumenting program code to reduce the 
% valid/allowed targets of indirect forwardedge transfers. We study the design
% of FineIBT on the x86-64 architecture, and implement and evaluate
% it on Linux and the LLVM
% toolchain. We designed FineIBT’s instrumentation to be compact,
% incurring low runtime and memory overheads, and generic, so as to
% support different CFI policies. Our prototype implementation incurs
% negligible runtime slowdowns (≈0%–1.94% in SPEC CPU2017 and
% ≈0%–1.92% in real-world applications) outperforming Clang-CFI.
% Lastly, we investigate the effectiveness/security and compatibility
% of FineIBT using the ConFIRM CFI benchmarking suite, demonstrating that our instrumentation provides complete coverage in
% the presence of modern software features, while supporting a wide
% range of CFI policies with the same, predictable performance


\project{Eclipse}{eclipse}
% NOT-OK: XXX
Memory corruption and speculative execution have often been studied separately,
yet attacks can combine them to defeat mitigations such as address-space
randomization and pointer authentication. Speculative probing, for example,
can test attacker-supplied pointers without causing the architectural fault
that would otherwise terminate the victim. My collaborators and I developed
Eclipse to block this class of speculative memory-error abuse. The compiler
identifies susceptible indirect branches and propagates artificial data
dependencies to their operands, preventing attacker-controlled pointers from
being used along mispredicted execution paths. This targets the attack's
speculative computation without requiring every side channel or memory error
to be eliminated, and without stopping speculation wholesale with fences.
We demonstrated protection against speculative-probing and PACMAN-style
attacks across userland and kernel settings on x86-64 and ARM. Our x86-64
performance evaluation measured approximately 0--9.5\% slowdown on SPEC
CPU2017, up to 8.6\% in real applications, and negligible overhead in the
Linux kernel, supporting the practicality of selective compiler hardening.

% Historically, researchers have treated memory safety-based and
% speculative execution attacks as two separate domains. Recent work
% has introduced Speculative Memory-error Abuse (SMA) attacks,
% which combine memory corruption vulnerabilities with Spectrelike primitives.
% Using SMA, an attacker can leak sensitive program
% information and defeat a wide variety of memory-corruption mitigations,
% including (K)ASLR, software-based XOM, and even ARM
% PA, eventually carrying out an end-to-end (architecturally-visible)
% exploit. We present Eclipse: a novel protection scheme against
% SMA attacks. Eclipse works by propagating artificial data dependencies onto
% sensitive data, preventing the CPU from using attackercontrolled data
% during speculative execution. We demonstrate that
% Eclipse provides comprehensive protection against speculativeprobing
% and PACMAN-style attacks, two prominent examples of
% SMA attacks that target both the x86(-64) and ARM architectures.
% We evaluate the performance of Eclipse on x86-64 and demonstrate that
% it introduces minimal overhead, compared to alternative
% hardening approaches, incurring ≈0%–9.5% slowdown on SPEC
% CPU 2017, up to 8.6% slowdown in real-world applications, and
% negligible overhead in the Linux kernel.


\project{EPF}{epf}
% NOT-OK: XXX
In EPF, we examined a blind spot in kernel defenses that isolate user-controlled
memory from kernel memory. BPF programs originate with users but reside in
kernel space, where the kernel legitimately accesses and executes them.
An attacker can therefore use their in-memory representation to encode exploit
payloads that remain accessible despite strong kernel-user isolation.
We developed two techniques, BPF-Reuse and BPF-ROP, and combined them with four
real-world vulnerabilities to demonstrate privilege escalation on 32- and
64-bit x86 systems. We then designed defenses that treat BPF as a separate
execution domain: isolate its instructions from ordinary kernel data, prevent
data from being interpreted as BPF code, and preserve the integrity of BPF
execution. Our Linux implementations incurred less than 3\% overhead on
BPF-heavy workloads. EPF thus contributes both an exploitation method and
concrete protections, showing why kernel hardening must account for
user-supplied code even after it has crossed into a trusted address space.

% The OS kernel is at the forefront of a system’s security. 
% Therefore, its own security is crucial for the correctness and
% integrity of user applications. With a plethora of bugs continuously 
% discovered in OS kernel code, defenses and mitigations are essential 
% for practical kernel security. One important
% defense strategy is to isolate user-controlled memory from
% kernel-accessible memory, in order to mitigate attacks like
% ret2usr and ret2dir. We present EPF (Evil Packet Filter): a new
% method for bypassing various (both deployed and proposed)
% kernel isolation techniques by abusing the BPF infrastructure
% of the Linux kernel: i.e., by leveraging BPF code, provided by
% unprivileged users/programs, as attack payloads.
% We demonstrate two different EPF instances, namely BPF-Reuse and
% BPF-ROP, which utilize malicious BPF payloads to mount
% privilege escalation attacks in both 32- and 64-bit x86 platforms.
% We also present the design, implementation, and evaluation of a set of
% defenses to enforce the isolation between BPF
% instructions and benign kernel data, and the integrity of BPF
% program execution, effectively providing protection against
% EPF-based attacks. Our implemented defenses show minimal
% overhead (< 3%) in BPF-heavy tasks.


\section{Robustness of ML Systems}
% NOT-OK: XXX
Intro about ml robustness, and the citations of pixeldp...
XXX Adversarial examples can make a neural network misclassify a slightly altered
input. With collaborators at Columbia, I helped develop PixelDP, which uses

\project{PixelDP}{pixeldp}
% NOT-OK: XXX
Adversarial examples can make a neural network misclassify a slightly altered
input. With collaborators at Columbia, I helped develop PixelDP, which uses
differential privacy to certify robustness against norm-bounded perturbations.
It adds calibrated noise to a model's computation and checks, at inference
time, whether an individual prediction remains stable within a specified
attack budget. The construction applies across model architectures; we
demonstrated certified $\ell_2$ robustness for Inception on ImageNet.
An autoencoder front end can also protect an existing model with limited
fine-tuning, providing a firewall-like deployment option. PixelDP extends
formal guarantees to large networks while exposing the tradeoff between
accuracy and the strength of the robustness bound.

\section{Vulnerability Remediation}
% NOT-OK: XXX
Research on vulnerability discovery matters most when teams can act on what is
found. My industry work addresses the operational problem of deciding which
known issues warrant attention first in a large production environment.

\project{Code-Aware Agent}{caa,caa_patent}
% NOT-OK: XXX
At Lacework, on Peter O'Hearn's Code Security team, I worked with Patrice
Godefroid on CAA. Package inventories and vulnerability databases can produce
more findings than a team can remediate quickly, while saying little about
whether an affected package runs in production. CAA continuously monitors
file-system activity, records when package files are opened or executed, and
sends package-activity information to the security platform. Correlating this
evidence with vulnerability findings lets teams focus remediation on packages
active in their environments. The monitoring must be inexpensive enough to
remain enabled throughout production execution. A publicly reported deployment
study covering 2,428 machines found that 79\% of the vulnerable packages were
inactive during the observation period. This result gives a concrete measure
of how runtime evidence can reduce the immediate remediation workload; the
activity classification continues to update as workloads change. The project
connects my research on program behavior to the engineering requirements of
an always-running security product.

\project{POSIX abstractions}{posix}
My earlier work asked whether the abstractions standardized by POSIX still
match the needs of modern applications. With collaborators at Columbia, I
built \texttt{libtrack} to combine detailed tracing of library calls with
static analysis at scale, including a study of 1.1 million Android
applications. Across Android, OS X, and Ubuntu consumer workloads, we found
that platform-specific frameworks increasingly provide abstractions for
graphics, inter-process communication, and structured storage that POSIX does
not directly support. Although these systems address similar needs, their
interfaces are diverging rather than converging on a shared standard.
Extension mechanisms such as \texttt{ioctl} conceal the semantics of these
operations from generic OS mediation, shifting responsibility to libraries
and drivers and complicating both protection and portability. The study
identified obsolete, unused, and missing abstractions, and argued for revising
the standard around measured application needs. This measurement-driven
approach remains central to how I want to study the interfaces required by
agentic workloads.


\section{Future Directions}
Kernel protection is a fertile and challenging research area that I plan to
continue working on, as kernel exploitation has currently turned into one of

\
\vspace{-0.75em}

{\noindent~\circled{1}~\textbf{Kernel Security and Transient Execution Attacks.}}
The second direction revisits the question raised by my POSIX study: which OS
abstractions does a new class of applications actually need? I plan to measure
agentic workloads as they create subprocesses, read and modify files, access
networks, and invoke external tools. Comparing those traces with the privileges
granted to each step can reveal where a missing abstraction forces developers
to use broad permissions or ad hoc wrappers. The result should be a
measurement-based account of agent behavior and concrete requirements for
interfaces and policies that preserve useful work while containing a
compromised or mistaken step. SysXCHG provides a starting point for asking
whether privileges can follow the work being performed more precisely.

\
\vspace{-0.75em}

{\noindent~\circled{1}~\textbf{Kernel Security and Transient Execution Attacks.}}
The second direction revisits the question raised by my POSIX study: which OS
abstractions does a new class of applications actually need? I plan to measure
agentic workloads as they create subprocesses, read and modify files, access
networks, and invoke external tools. Comparing those traces with the privileges
granted to each step can reveal where a missing abstraction forces developers
to use broad permissions or ad hoc wrappers. The result should be a
measurement-based account of agent behavior and concrete requirements for
interfaces and policies that preserve useful work while containing a
compromised or mistaken step. SysXCHG provides a starting point for asking
whether privileges can follow the work being performed more precisely.

\
\vspace{-0.75em}

{\noindent~\circled{3}~\textbf{Diversification and Debloating.}}
Diversifying the internal structure of system components, coupled with the
The third direction treats the agent harness itself as a stateful system.
Building on RESTler's dependency inference and Pythia's feedback-guided
mutations, I want to generate multi-step tool interactions that vary
permissions, environment state, and external responses. Active checkers can
test properties such as authorization, isolation between tasks, and recovery
after a failed tool call, including failures that do not produce a crash. I
would evaluate these methods on real harnesses, with each finding reduced to
a reproducible trace and a clear violated property. Together, these three
directions would let testing evidence guide targeted hardening and
remediation. I would welcome the opportunity to pursue this agenda with
Bloomberg and help build dependable software systems at scale.

renewcommand{\refname}{References}
\small{
\bibliographystyle{abbrv}
\bibliography{statement}
}

\end{document}
